\documentclass[12pt]{article}

\usepackage{amsmath}
\usepackage[margin=1in]{geometry}
\usepackage{fancyhdr}
\usepackage{amsthm}
\newtheorem{lemma}{Lemma}
\usepackage{braket}
\usepackage{amssymb}
\usepackage{pgfplots}
\usepackage{tikz}
\usepackage{bbm}
\usepackage{bm}
\usetikzlibrary{shapes.geometric}
\usepackage{enumerate}
\usepackage[shortlabels]{enumitem}


\pagestyle{fancy}
\fancyhead[l]{Hudson Benites}
\fancyhead[c]{Day 2}
\fancyfoot[c]{\thepage}
\renewcommand{\headrulewidth}{0.2pt}
\setlength{\headheight}{15pt}
\fancyhead[r]{\today}

\newcommand{\magnitude}[1]{\left\lvert #1 \right\rvert}

\begin{document}
REPLACE k s with KAPPA
\begin{enumerate}[label=\Alph*., leftmargin=*]
    \item $\kappa =\frac{\bm{v}\times \bm{a}}{{\vert \bm{v}\vert}^3}.$ \begin{proof}
      Consider the following setup
      \begin{center}
      \begin{tikzpicture}
        \draw[->] (0,0) -- (2,2) node[midway, above left]{$\bm{a}$};
        \draw[->] (0,0) -- (0, 2) node[midway, above left]{$\bm{a^{\perp}}$};
        \draw[->] (0,0) -- (4,0) node[pos=1, right]{$\bm{v}$};
        \draw[->] (0,0) -- (2,0) node[midway, below]{$\bm{a}^{\parallel}$};
        \draw (0.6,0) arc (0:45:0.6) node[midway, right]{$\theta$};
      \end{tikzpicture}
    \end{center}
    From the above picture, we can see that $\sin(\theta)=\frac{\magnitude{\bm{a}^{\perp}}}{\magnitude{\bm{a}}}$, which implies that \begin{equation*}
      \magnitude{\bm{a}}\sin(\theta)=\magnitude{\bm{\alpha}^{\perp}}.
    \end{equation*}
     Note that $\magnitude{\bm{v}\times\bm{a}}=\magnitude{\bm{v}}\magnitude{\bm{a}}\sin(\theta)$. Using the equation from above, we see that $\magnitude{\bm{v}\times\bm{a}}=\magnitude{\bm{v}}\magnitude{\bm{a}^{\perp}}$. Dividing by $\magnitude{\bm{v}}^3$, we compute\begin{equation*}
      \frac{\magnitude{\bm{v}\times\bm{a}}}{\magnitude{\bm{v}}^3}=\frac{\magnitude{\bm{v}}\magnitude{\bm{a}^{\perp}}}{\magnitude{\bm{v}}^3}=\frac{\magnitude{\bm{a}}^{\perp}}{\magnitude{\bm{v}}^2}=\kappa.
     \end{equation*}
     We are done.
    \end{proof}
    \item $\kappa = \frac{\dot{x}\ddot{y}-\ddot{x}\dot{y}}{{(\dot{x}^2+\dot{y}^2)}^{\frac{3}{2}}}$.\begin{proof}
    We can embed $\bm{\alpha}$ into $\mathbb{R}^3$ by constructing an equivalent $\bm{\alpha}: \mathbb{R}\rightarrow\mathbb{R}^3$, $\bm{\alpha}=(x(u),y(u),0)$. We can see that $\bm{\dot{\alpha}}=(\dot{x}(u),\dot{y}(u),0)$ and $\bm{\ddot{\alpha}}=(\ddot{x}(u),\ddot{y}(u),0)$. 
    We compute that\begin{equation*}
      \bm{\dot{\alpha}}\times \bm{\ddot{\alpha}}=\dot{x}\ddot{y}-\ddot{x}\dot{y}.
    \end{equation*}
    We also compute that\begin{equation*}
      \magnitude{\bm{\dot{\alpha}}}=\sqrt{\dot{x}^2+\dot{y}^2}.
    \end{equation*}
    Using the identity proven in (A), we see that\begin{equation*}
      \kappa = \frac{ \bm{\dot{\alpha}}\times \bm{\ddot{\alpha}}}{(\magnitude{\bm{\dot{\alpha}}})^3}=\frac{\dot{x}\ddot{y}-\ddot{x}\dot{y}}{(\dot{x}^2+\dot{y}^2)^{3/2}}
    \end{equation*}
    Thus, we are done.
      \end{proof}
    \item $\frac{d\bm{T}}{ds}=\kappa N$. \begin{proof} $\bm{T}$ exists only for a regular curve, and therefore we assume $\bm{\alpha}$ to be regular. By definition, $\bm{T}=\frac{\bm{\dot{\alpha}}}{\vert\bm{\dot{\alpha}}\vert}$. Because $\bm{\alpha}$ is regular
          there exists an arc length parameterization $s(u)$ that gives a $\bm{T}=\bm{\dot{\alpha}}(s)$, because $\vert \bm{\dot{\alpha}}(s)\vert=1$ for all $u$. Differentiating with respect to $s$, we compute \begin{equation*}
            \frac{d\bm{T}}{ds}=\bm{\ddot{\alpha}}(s).
          \end{equation*} We now show that $a^{\perp}=a$ for all $u$. Note that because $\bm{\alpha}$ is arc length parameterized, $\vert \bm{\dot{\alpha}} \vert = 1$ always, and thus $1=\bm{\dot{\alpha}}\cdot\bm{\dot{\alpha}}$. We compute
          \begin{equation*}
            \frac{d}{du}(1)=\frac{d}{du}(\bm{\dot{\alpha}}\cdot\bm{\dot{\alpha}}),
          \end{equation*} and so
          \begin{equation*}
            0 = \bm{\ddot{\alpha}}\cdot\bm{\dot{\alpha}}+\bm{\dot{\alpha}}\cdot\bm{\ddot{\alpha}}=2(\bm{\dot{\alpha}}\cdot\bm{\ddot{\alpha}}).
          \end{equation*} The above equation implies that $\bm{\dot{\alpha}}$ is perpendicular to $\bm{\ddot{\alpha}}$, since their dot product is 0.
          Therefore, there is no component of acceleration that points in the direction $v$; as such, it must be the case that $\bm{\ddot{\alpha}}^{\perp}=\bm{\ddot{\alpha}}$.
          Note that $\kappa\bm{N} = \frac{\vert \bm{\ddot{\alpha}}^{\perp}\vert}{{\vert \bm{\dot{\alpha}}\vert}^2}\frac{\bm{\ddot{\alpha}}^{\perp}}{\vert\bm{\ddot{\alpha}}^{\perp}\vert}$. Since $\bm{\alpha}$ is arc length parameterized,
          $\vert \bm{\dot{\alpha}}\vert=1$ always, so $\kappa N=\bm{\ddot{\alpha}}^{\perp}$, which we have shown to be equal to $\bm{\ddot{\alpha}}$. Therefore, $\frac{d\bm{T}}{ds}=\kappa N$, and we are done. 
        \end{proof}
      \item By assumption $\bm{\alpha}$ is regular, and therefore there exists an arc length parameterization (A.L.P) of $\bm{\alpha}$. Let $\bm{\alpha}$ be A.L.P-ed. Furthermore, let \begin{equation*} 
        \bm{\beta}(v)=\frac{1}{\kappa(u_0)}(\bm{N}(u_0))(1-\cos(v))+\frac{1}{\kappa(u_0)}\sin(v)\bm{T}(u_0)+\alpha(u_0).
      \end{equation*} Note that $\cos(v)$ and $\sin(v)$ are infinitely differentiable, and thus $\bm{\beta}$ is smooth. Furthermore, note that $\bm{\dot{\beta}}=\frac{1}{\kappa(u_0)}\sin(v)\bm{N}(u_0)+\frac{1}{\kappa(u_0)}\cos(v)\bm{T}(u_0)$, which can never be the zero vector.
            Therefore, $\bm{\beta}$ is regular, and and A.L.P exists. Let $\bm{\beta}$ be A.L.P-ed. We show that $\beta$ satisifes (1)-(4).
            \begin{enumerate}[label=(\arabic*)]
              \item We compute\begin{equation*}
                \bm{\beta}(0)=\frac{1}{\kappa(u_0)}(-\cos(0)+1)+0+\bm{\alpha}(u_0)=\bm{\alpha}(u_0).
              \end{equation*}
              \item We compute 
              \begin{equation*}
                \bm{T}_{\beta}(0)=\frac{\bm{\dot{\beta}}(0)}{\magnitude{\bm{\dot{\beta}}(0)}}=\frac{\sin(0)\bm{N}(u_0)+\cos(0)\bm{T}(u_0)}{\frac{1}{\kappa(u_0)}\kappa(u_0)}.
              \end{equation*} Therefore, $\bm{T}_{\bm{\beta}}(0)=\bm{T}(u_0)$.
              \item We compute \begin{equation*}
                \bm{N}_{\bm{\beta}}(0)=\frac{\bm{\ddot{\beta}}^{\perp}(0)}{\magnitude{\bm{\ddot{\beta}}^{\perp}(0)}}=\frac{\cos(0)\bm{N}(u_0)-\sin(0)\bm{T}(u_0)}{\frac{1}{\kappa(u_0)}\kappa(u_0)}.
              \end{equation*} Therefore, $\bm{N}_{\bm{\beta}}(0)=\bm{N}(u_0)$. From (C), we know that because $\bm{\beta}$ is A.L.P-ed, $\bm{\ddot{\beta}}^{\perp}=\bm{\ddot{\beta}}$, and therefore the equality above holds.
              \item We compute\begin{equation*}
                  \kappa_{\bm{\beta}}(0)=\frac{\magnitude{\bm{\ddot{\beta}}(0)}}{\magnitude{\bm{\dot{\beta}}(0)}}=\frac{\frac{1}{\kappa(u_0)}}{(\frac{1}{\kappa(u_0)})^2}=\kappa(u_0).
              \end{equation*}
              Therefore, $\kappa_{\bm{\beta}}=\kappa(u_0)$.
            \end{enumerate}
    \end{enumerate}
\end{document}